% ============================================================
\chapter{Th\'eor\`eme de Carath\'eodory}
\label{ch:caratheodory}
% ============================================================

\section{Introduction}

Comment construire une mesure~? La question peut sembler circulaire~: on veut mesurer des ensembles, mais pour d\'efinir rigoureusement ce que signifie \og mesurer \fg, il faut d'abord disposer d'une th\'eorie. C'est Constantin Carath\'eodory qui, en 1914, a tranch\'e ce n\oe ud gordien avec une construction d'une \'el\'egance remarquable. Son id\'ee~: partir d'une notion plus grossi\`ere, la \emph{mesure ext\'erieure}, qui a le d\'efaut de ne pas \^etre additive en g\'en\'eral, puis s\'electionner parmi tous les ensembles ceux sur lesquels cette mesure ext\'erieure se comporte bien. Les ensembles \og bien mesurables \fg au sens de Carath\'eodory forment automatiquement une $\sigma$-alg\`ebre, et la mesure ext\'erieure restreinte \`a cette $\sigma$-alg\`ebre est une vraie mesure.

Cette construction, aussi \'el\'egante qu'abstraite, est la machine universelle qui produit la mesure de Lebesgue, la mesure de Hausdorff, et bien d'autres. C'est la pierre angulaire sur laquelle repose tout l'\'edifice de l'int\'egration moderne.

\section{Mesures ext\'erieures}

\begin{definition}[Mesure ext\'erieure]
Soit $X$ un ensemble. Une \emph{mesure ext\'erieure} sur $X$ est une application
$\mu^* : \mathcal{P}(X) \to [0, +\infty]$ v\'erifiant~:
\begin{enumerate}
    \item[(i)] $\mu^*(\varnothing) = 0$\,;
    \item[(ii)] (monotonie) si $A \subset B$, alors $\mu^*(A) \leq \mu^*(B)$\,;
    \item[(iii)] ($\sigma$-sous-additivit\'e) pour toute suite $(A_n)_{n \in \N}
    \subset \mathcal{P}(X)$~:
    \[
    \mu^*\!\left(\bigcup_{n=0}^{\infty} A_n\right) \leq \sum_{n=0}^{\infty}
    \mu^*(A_n).
    \]
\end{enumerate}
\end{definition}

\begin{remarque}
Une mesure ext\'erieure est d\'efinie sur \emph{tous} les sous-ensembles de $X$, mais
elle n'est en g\'en\'eral pas $\sigma$-additive. L'id\'ee de Carath\'eodory est
d'identifier les ensembles sur lesquels elle est additive.
\end{remarque}

\section{Ensembles Carath\'eodory-mesurables}

\begin{definition}[Mesurabilit\'e au sens de Carath\'eodory]
\label{def:caratheodory_mesurable}
Soit $\mu^*$ une mesure ext\'erieure sur $X$. Un ensemble $A \subset X$ est dit
\emph{$\mu^*$-mesurable} (ou \emph{Carath\'eodory-mesurable}) si~:
\[
\forall\, E \subset X, \quad \mu^*(E) = \mu^*(E \cap A) + \mu^*(E \cap A^c).
\]
On note $\mathcal{M}^*$ la collection de tous les ensembles $\mu^*$-mesurables.
\end{definition}

\begin{intuition}[title=\textbf{Interpr\'etation du crit\`ere de Carath\'eodory}]
La condition dit que $A$ \og d\'ecoupe \fg proprement tout ensemble $E$~: la mesure
ext\'erieure de $E$ est exactement la somme des mesures ext\'erieures des deux
morceaux $E \cap A$ et $E \cap A^c$. Par la sous-additivit\'e, l'in\'egalit\'e
$\mu^*(E) \leq \mu^*(E \cap A) + \mu^*(E \cap A^c)$ est toujours vraie. Il suffit
donc de v\'erifier l'in\'egalit\'e inverse~:
$\mu^*(E) \geq \mu^*(E \cap A) + \mu^*(E \cap A^c)$.
\end{intuition}

\section{Le th\'eor\`eme principal}

\begin{theoreme}[Carath\'eodory]
\label{thm:caratheodory}
Soit $\mu^*$ une mesure ext\'erieure sur $X$. Alors~:
\begin{enumerate}
    \item $\mathcal{M}^*$ est une $\sigma$-alg\`ebre\,;
    \item la restriction $\mu = \mu^*|_{\mathcal{M}^*}$ est une mesure compl\`ete
    sur $(X, \mathcal{M}^*)$.
\end{enumerate}
\end{theoreme}

\begin{proof}
La preuve se d\'ecompose en plusieurs \'etapes.

\emph{\'Etape 1~: $\mathcal{M}^*$ est une alg\`ebre.}

Clairement $X \in \mathcal{M}^*$ (prendre $A = X$~: $\mu^*(E) = \mu^*(E) +
\mu^*(\varnothing) = \mu^*(E)$). Si $A \in \mathcal{M}^*$, alors $A^c \in
\mathcal{M}^*$ car la condition de Carath\'eodory est sym\'etrique en $A$ et $A^c$.

Montrons la stabilit\'e par r\'eunion finie. Soient $A, B \in \mathcal{M}^*$. Pour
tout $E \subset X$~:
\begin{align*}
\mu^*(E) &= \mu^*(E \cap A) + \mu^*(E \cap A^c) \\
&= \mu^*(E \cap A \cap B) + \mu^*(E \cap A \cap B^c) + \mu^*(E \cap A^c \cap B)
+ \mu^*(E \cap A^c \cap B^c).
\end{align*}
Or $E \cap (A \cup B) = (E \cap A \cap B) \cup (E \cap A \cap B^c) \cup
(E \cap A^c \cap B)$ et $E \cap (A \cup B)^c = E \cap A^c \cap B^c$. Par
sous-additivit\'e~:
\[
\mu^*(E \cap (A \cup B)) \leq \mu^*(E \cap A \cap B) + \mu^*(E \cap A \cap B^c)
+ \mu^*(E \cap A^c \cap B).
\]
Donc $\mu^*(E) \geq \mu^*(E \cap (A \cup B)) + \mu^*(E \cap (A \cup B)^c)$, et
l'in\'egalit\'e inverse vient de la sous-additivit\'e. Ainsi $A \cup B \in
\mathcal{M}^*$.

\emph{\'Etape 2~: additivit\'e finie.}

Si $A, B \in \mathcal{M}^*$ sont disjoints, en prenant $E = A \cup B$~:
$\mu^*(A \cup B) = \mu^*((A \cup B) \cap A) + \mu^*((A \cup B) \cap A^c)
= \mu^*(A) + \mu^*(B)$.

Par r\'ecurrence, pour $A_1, \ldots, A_n \in \mathcal{M}^*$ deux \`a deux disjoints~:
$\mu^*\!\left(\bigsqcup_{k=1}^n A_k\right) = \sum_{k=1}^n \mu^*(A_k)$.

\emph{\'Etape 3~: $\sigma$-additivit\'e et stabilit\'e par r\'eunion d\'enombrable.}

Soit $(A_n)_{n \geq 1} \subset \mathcal{M}^*$ une suite disjointe. Posons
$S_n = \bigsqcup_{k=1}^n A_k$ et $A = \bigsqcup_{k=1}^\infty A_k$. Par l'\'etape 1,
$S_n \in \mathcal{M}^*$, donc pour tout $E$~:
\[
\mu^*(E) = \mu^*(E \cap S_n) + \mu^*(E \cap S_n^c)
\geq \sum_{k=1}^n \mu^*(E \cap A_k) + \mu^*(E \cap A^c)
\]
car $S_n^c \supset A^c$. En faisant $n \to \infty$~:
\[
\mu^*(E) \geq \sum_{k=1}^{\infty} \mu^*(E \cap A_k) + \mu^*(E \cap A^c)
\geq \mu^*(E \cap A) + \mu^*(E \cap A^c)
\]
par sous-additivit\'e de $\mu^*$. L'in\'egalit\'e inverse \'etant toujours vraie,
$A \in \mathcal{M}^*$.

En prenant $E = A$~: $\mu^*(A) = \sum_{k=1}^{\infty} \mu^*(A_k) + \mu^*(\varnothing)
= \sum_k \mu^*(A_k)$, ce qui donne la $\sigma$-additivit\'e.

\emph{\'Etape 4~: compl\'etude.}

Soit $N \in \mathcal{M}^*$ avec $\mu^*(N) = 0$ et $A \subset N$. Pour tout
$E \subset X$~: $\mu^*(E \cap A) \leq \mu^*(N) = 0$ et $\mu^*(E \cap A^c) \leq
\mu^*(E)$, donc $\mu^*(E \cap A) + \mu^*(E \cap A^c) \leq \mu^*(E)$. L'in\'egalit\'e
inverse est la sous-additivit\'e, donc $A \in \mathcal{M}^*$.
\end{proof}

\section{Application aux pr\'e-mesures}

\begin{theoreme}[Extension de Carath\'eodory]
\label{thm:extension_caratheodory}
Soit $\mu_0$ une pr\'e-mesure sur une alg\`ebre $\mathcal{A}_0 \subset
\mathcal{P}(X)$. D\'efinissons~:
\[
\mu^*(A) = \inf\left\{\sum_{n=0}^{\infty} \mu_0(A_n) :
(A_n) \subset \mathcal{A}_0,\, A \subset \bigcup_{n=0}^{\infty} A_n\right\},
\quad A \subset X.
\]
Alors~:
\begin{enumerate}
    \item $\mu^*$ est une mesure ext\'erieure sur $X$.
    \item $\mathcal{A}_0 \subset \mathcal{M}^*$ (tout \'el\'ement de l'alg\`ebre est
    Carath\'eodory-mesurable).
    \item $\mu^*|_{\mathcal{A}_0} = \mu_0$ (l'extension co\"incide avec la
    pr\'e-mesure sur l'alg\`ebre).
    \item En particulier, $\mu = \mu^*|_{\sigma(\mathcal{A}_0)}$ est une mesure
    sur $\sigma(\mathcal{A}_0)$ qui \'etend $\mu_0$.
\end{enumerate}
\end{theoreme}

\begin{proof}
\emph{Point 1.} $\mu^*(\varnothing) = 0$ (prendre $A_n = \varnothing$). La
monotonie est claire. Pour la $\sigma$-sous-additivit\'e, soit $(B_k)$ une suite et
$\varepsilon > 0$. Pour chaque $k$, il existe $(A_n^{(k)}) \subset \mathcal{A}_0$
avec $B_k \subset \bigcup_n A_n^{(k)}$ et $\sum_n \mu_0(A_n^{(k)}) \leq
\mu^*(B_k) + \varepsilon/2^{k+1}$. Alors $\bigcup_k B_k \subset
\bigcup_{k,n} A_n^{(k)}$ et
\[
\mu^*\!\left(\bigcup_k B_k\right) \leq \sum_{k,n} \mu_0(A_n^{(k)}) \leq
\sum_k \mu^*(B_k) + \varepsilon.
\]
En faisant $\varepsilon \to 0$, on obtient la $\sigma$-sous-additivit\'e.

\emph{Point 2.} Soit $A \in \mathcal{A}_0$ et $E \subset X$. Il suffit de montrer
$\mu^*(E) \geq \mu^*(E \cap A) + \mu^*(E \cap A^c)$. Soit $\varepsilon > 0$ et
$(A_n) \subset \mathcal{A}_0$ avec $E \subset \bigcup_n A_n$ et $\sum_n \mu_0(A_n)
\leq \mu^*(E) + \varepsilon$. Comme $\mathcal{A}_0$ est une alg\`ebre,
$A_n \cap A$ et $A_n \cap A^c$ sont dans $\mathcal{A}_0$ et
$\mu_0(A_n) = \mu_0(A_n \cap A) + \mu_0(A_n \cap A^c)$ (additivit\'e finie). Donc~:
\begin{align*}
\mu^*(E) + \varepsilon &\geq \sum_n \mu_0(A_n \cap A) + \sum_n \mu_0(A_n \cap A^c)
\\ &\geq \mu^*(E \cap A) + \mu^*(E \cap A^c).
\end{align*}

\emph{Point 3.} Pour $A \in \mathcal{A}_0$, $\mu^*(A) \leq \mu_0(A)$ (prendre
$A_0 = A$, $A_n = \varnothing$). Pour l'in\'egalit\'e inverse, soit $(A_n) \subset
\mathcal{A}_0$ avec $A \subset \bigcup_n A_n$. Posons $B_N = A \cap
\bigl(\bigcup_{n=0}^N A_n\bigr) \in \mathcal{A}_0$. Alors $B_N \uparrow A$ et
$B_N \subset \bigcup_{n=0}^N A_n$. Par sous-additivit\'e finie et la pr\'e-mesure~:
$\mu_0(B_N) \leq \sum_{n=0}^N \mu_0(A_n)$. Par continuit\'e croissante de la
pr\'e-mesure (proposition~\ref{prop:critere_sigma_add})~:
$\mu_0(A) = \lim_N \mu_0(B_N) \leq \sum_n \mu_0(A_n)$. En prenant l'infimum~:
$\mu_0(A) \leq \mu^*(A)$.
\end{proof}

\section{Unicit\'e de l'extension}

\begin{theoreme}[Unicit\'e]
\label{thm:unicite_extension}
Si la pr\'e-mesure $\mu_0$ est $\sigma$-finie sur $\mathcal{A}_0$, alors l'extension
$\mu$ \`a $\sigma(\mathcal{A}_0)$ est unique.
\end{theoreme}

\begin{proof}
Par le th\'eor\`eme d'unicit\'e des mesures (th\'eor\`eme~\ref{thm:unicite_mesures}),
appliqu\'e au $\pi$-syst\`eme $\mathcal{A}_0$ (une alg\`ebre est un $\pi$-syst\`eme).
La $\sigma$-finitude de $\mu_0$ fournit la suite exhaustive requise.
\end{proof}

\begin{formules}[title=\textbf{R\'esum\'e de la construction de Carath\'eodory}]
\[
\text{Pr\'e-mesure } \mu_0 \text{ sur } \mathcal{A}_0
\xrightarrow{\text{mesure ext.}} \mu^* \text{ sur } \mathcal{P}(X)
\xrightarrow{\text{restriction}} \mu \text{ sur } \mathcal{M}^*
\supset \sigma(\mathcal{A}_0).
\]
\end{formules}

\section{Crit\`ere de Carath\'eodory pour les mesures m\'etriques}

\begin{definition}[Mesure ext\'erieure m\'etrique]
Soit $(X, d)$ un espace m\'etrique. Une mesure ext\'erieure $\mu^*$ sur $X$ est dite
\emph{m\'etrique} si~:
\[
d(A, B) > 0 \implies \mu^*(A \cup B) = \mu^*(A) + \mu^*(B),
\]
o\`u $d(A, B) = \inf\{d(a, b) : a \in A,\, b \in B\}$.
\end{definition}

\begin{theoreme}
Si $\mu^*$ est une mesure ext\'erieure m\'etrique sur $(X, d)$, alors tout ensemble
bor\'elien est $\mu^*$-mesurable~: $\mathcal{B}(X) \subset \mathcal{M}^*$.
\end{theoreme}

\begin{proof}
Il suffit de montrer que tout ferm\'e $F$ est $\mu^*$-mesurable. Soit $E \subset X$
et posons $E_n = \{x \in E \cap F^c : d(x, F) \geq 1/n\}$. Alors $E_n \uparrow
E \cap F^c$ et $d(E \cap F, E_n) \geq 1/n > 0$, donc
$\mu^*(E \cap F \cup E_n) = \mu^*(E \cap F) + \mu^*(E_n)$ par la propri\'et\'e
m\'etrique.

Par monotonie, $\mu^*(E) \geq \mu^*(E \cap F \cup E_n) = \mu^*(E \cap F) +
\mu^*(E_n)$. Il reste \`a montrer que $\mu^*(E_n) \to \mu^*(E \cap F^c)$.

Posons $D_n = E_{n+1} \setminus E_n$. Pour $n$ et $m$ avec $m \geq n + 2$,
$d(D_n, D_m) > 0$, donc les ensembles $D_n$ de parit\'e fix\'ee sont \`a distance
positive deux \`a deux. Par la propri\'et\'e m\'etrique et la $\sigma$-sous-additivit\'e~:
\[
\mu^*(E \cap F^c) \leq \mu^*(E_n) + \sum_{k=n}^{\infty} \mu^*(D_k).
\]
Si $\sum_k \mu^*(D_k) < \infty$, le reste tend vers z\'ero et
$\mu^*(E_n) \to \mu^*(E \cap F^c)$. Si la s\'erie diverge, on montre que
$\mu^*(E) = +\infty$ et l'in\'egalit\'e est triviale.
\end{proof}

\section{Application~: mesure de Hausdorff}

\begin{definition}[Mesure de Hausdorff]
Soit $(X, d)$ un espace m\'etrique et $s \geq 0$. Pour $\delta > 0$, on d\'efinit~:
\[
\mathcal{H}^s_\delta(A) = \inf\left\{\sum_{n=1}^{\infty} (\operatorname{diam}
U_n)^s : A \subset \bigcup_{n=1}^{\infty} U_n,\,
\operatorname{diam} U_n \leq \delta\right\}.
\]
La \emph{mesure de Hausdorff $s$-dimensionnelle} est~:
\[
\mathcal{H}^s(A) = \lim_{\delta \to 0} \mathcal{H}^s_\delta(A) =
\sup_{\delta > 0} \mathcal{H}^s_\delta(A).
\]
\end{definition}

\begin{proposition}
$\mathcal{H}^s$ est une mesure ext\'erieure m\'etrique. En cons\'equence, c'est une
mesure (bor\'elienne) sur $(X, \mathcal{B}(X))$.
\end{proposition}

\section{Exercices}

\begin{exercice}
V\'erifier en d\'etail que l'ensemble $\mathcal{M}^*$ du th\'eor\`eme de Carath\'eodory
est stable par intersection d\'enombrable.
\end{exercice}

\begin{exercice}
Soit $\mu^*$ une mesure ext\'erieure sur $X$. Montrer que si $\mu^*(A) = 0$, alors
$A \in \mathcal{M}^*$.
\end{exercice}

\begin{exercice}
Montrer que la mesure de Hausdorff $\mathcal{H}^n$ sur $\R^n$ est (\'a une constante
multiplicative pr\`es) la mesure de Lebesgue.
\end{exercice}

\begin{exercice}
Soit $\mu_0$ une pr\'e-mesure \emph{finie} sur une alg\`ebre $\mathcal{A}_0$.
Montrer que pour tout $A \in \sigma(\mathcal{A}_0)$ et tout $\varepsilon > 0$,
il existe $B \in \mathcal{A}_0$ tel que $\mu(A \triangle B) < \varepsilon$, o\`u
$A \triangle B = (A \setminus B) \cup (B \setminus A)$.
\end{exercice}

\begin{exercice}
Soit $X$ un ensemble non d\'enombrable et
$\mu^*(A) = \begin{cases} 0 & \text{si } A \text{ est d\'enombrable,} \\
1 & \text{sinon.} \end{cases}$
Montrer que $\mu^*$ est une mesure ext\'erieure et d\'eterminer $\mathcal{M}^*$.
\end{exercice}
